\documentclass[12pt,twoside]{article}
\usepackage[ngerman]{babel} 

\usepackage[T1]{fontenc}
\usepackage{lmodern}
\usepackage{graphicx}
\usepackage{amsmath,amssymb,amsfonts,amsthm}
\begin{document}
\textbf{Fragen:}
\begin{enumerate}
\item Ist es richtig, dass bei uns $T_p(E)$ also $T_p(L^2(\Omega,H))=1$ und $p=2$ gilt, da wir einen Hilbertraum haben?
\item Woher kommt genau die Wahl am Ende von $a_l^*$ und was ist $a_1^*$?
\item Wie bekomme ich die Konstanten $C_2$ und $q_2$?\\
Ich habe versucht mit Proposition 4.1 zu arbeiten. Danach gilt ja
\begin{align*}
\|D_l-\mathbb{E}(D_l)\|^2=\frac{1}{\pi m_l}(1-\frac{2}{\pi m_l})
\end{align*}
Also
\begin{align*}
\|D_l-\mathbb{E}(D_l)\|\leq\frac{1}{\pi m_l}(1-\frac{2}{\pi m_l})< \frac{1}{\pi}(\frac{1}{2})^l
\end{align*}
Also k"onnte man h wieder als 1/2 $q_2=1$ und $C_2=\frac{1}{\pi}$. Aber das ist doch zu ungenau. Ich habe auch hier lange versucht auf was anderes zu kommen.
\item Wie komme ich auf die Konstanten $C_3$ und $r$, die wenn ich richtig verstanden habe, wichtig f"ur das Optimierungsproblem sind?
\item Wie l"ose ich die neuen Optimierungsproblem, welches ich durch das Branchen bekomme?\\
Meine Idee ist immer wieder das Lemma 5.3 anzuwenden. Jedoch ist mir nicht genau klar wie ich das machen kann.\\
Meine erste Idee war, nachdem man z.B die Fallunterscheidung f"ur eine Komplente gemacht hat die Zweige des Baumes nicht mit $\leq$ und $\geq$ zu beschriften sondern wirklich direkt die Gleichheit anzunehmem und nur noch "uber die restlichen Komponenten zu minimieren. Die Nebenbedingung passt man dann an, indem man das $\delta^p$ neu bestimmen muss. Jedoch find ich das nicht gut.
Die andere kann ich gerade schwer so in Worten formulieren. \\
Ich hab mir nat"urlich auch "uberlegt es anders, als mit dem Lemma zu l"osen aber bin auf nichts gekommen.
\end{enumerate}
\end{document}